\documentclass[10pt,a4paper]{amsart}
\usepackage[margin=19mm]{geometry}
\usepackage{amsmath,amssymb,mathtools}
\usepackage[scr=rsfs,cal=euler]{mathalfa}
\usepackage{xfrac}
\usepackage[unicode,hidelinks,bookmarks=false]{hyperref}
\newcommand{\ie}{i.e.\ }
\newcommand{\cf}{cf.\ }
\newcommand{\resp}{respectively\ }
\newcommand{\Subsection}{\subsection}
\providecommand{\nobreakdash}{\nobreak\hbox{-}}
\setlength{\emergencystretch}{2em}
\multlinegap0pt
\usepackage{bm}
\usepackage{bbm}
\usepackage{dsfont}
\usepackage{xspace}
\usepackage{cancel}
\usepackage{mathtools}
\usepackage{amsthm}
\makeatletter
\@ifundefined{remark}{\newtheorem{remark}{Remark}}{}
\@ifundefined{theorem}{\newtheorem{theorem}{Theorem}}{}
\@ifundefined{lemma}{\newtheorem{lemma}[theorem]{Lemma}}{}
\makeatother
\newcommand*{\balpha}{\boldsymbol{\alpha}}
\newcommand*{\bbeta}{\boldsymbol{\beta}}
\newcommand*{\bdelta}{\boldsymbol{\delta}}
\newcommand*{\bgamma}{\boldsymbol{\gamma}}
\newcommand*{\botimes}{\boldsymbol{\otimes}}
\newcommand*{\cendot}{\boldsymbol{\cdot}}
\newcommand*{\df}{\mathfrak d}
\newcommand*{\id}{{\mathrm{id}}}
\title[The noncommutative KP hierarchy and its solution via descent]{The noncommutative KP hierarchy and its solution via descent algebra}
\author{Gordon Blower, Simon J.A. Malham}
\date{Source: arXiv:2510.01352v1 (CC BY 4.0)}
\begin{document}
\maketitle

\noindent\emph{Source and licence.} The noncommutative KP hierarchy and its solution via descent algebra. Version arXiv:2510.01352v1 (CC BY 4.0), distributed under the Creative Commons Attribution 4.0 International licence, as recorded in the arXiv licence metadata for this exact version and shown on the abstract page https://arxiv.org/abs/2510.01352v1. Full rights notice: licenses/arxiv-2510.01352-CC-BY-4.0.txt.\par\smallskip

\noindent\textit{Editorial scope.} Counted unit: the Theorem whose source label is \texttt{thm:Satocoeffsdescents} in the article (source0.tex lines 2889--2895, source label \texttt{thm:Satocoeffsdescents}), delivered together with the article's own complete original proof of it (source0.tex lines 2897--3069, one \texttt{proof} environment of 173 lines). No mathematics is written, repaired or re-derived: statement and proof are the article's own text line for line. Required context retained uncounted: thm:weightformula (lines 2547--2557); lemma:Satocoeffsviadescents (lines 2719--2724); eq:siglength3 (lines 2869--2877); rmk:mainthmstrategy (lines 2881--2887). Every definition, display and label of those lines is retained unchanged. Omitted from this package and not redistributed: 1--2546, 2558--2718, 2725--2868, 2878--2880, 2888--2888, 2896--2896, 3070--3528. Nothing inside a retained excerpt is omitted or altered: every retained body is delivered exactly as the article's lines are written, so no omission receipt of any kind accompanies this package. Citations: the retained excerpts carry no citation command at all, so no bibliography entry of the article is transcribed and no reference is redistributed. Reference adaptation: none was needed and none was made; every reference inside the delivered unit resolves inside it, so no unresolved reference or citation is produced. Printed slips kept literal: no printed word, symbol, hypothesis or number of the counted statement or of its proof was altered. Layout and class: the article's own class is \texttt{elsarticle} with packages including amsmath, amssymb, amsfonts, mathtools, forest, stackengine, relsize, setspace, xy; the wrapper is the lane's pinned \texttt{amsart} editorial wrapper at 10pt on A4, which restates the theorem-like environments the retained text needs and adds the article's own macro definitions that the retained text needs (\texttt{\textbackslash balpha}, \texttt{\textbackslash bbeta}, \texttt{\textbackslash bdelta}, \texttt{\textbackslash bgamma}, \texttt{\textbackslash botimes}, \texttt{\textbackslash cendot}, \texttt{\textbackslash df}, \texttt{\textbackslash id}), with extra wrapper packages (bm, bbm, dsfont, xspace, cancel, mathtools). The delivered numbering is the unit's own. Assets: no retained span contains a figure environment, a graphics-inclusion command, a PDF-page inclusion, a table or a TikZ picture; no image file is copied into this package, the package declares no asset dependency and no image file is redistributed with it. Licence: version arXiv:2510.01352v1 of this article is distributed under the Creative Commons Attribution 4.0 International licence, as recorded in the arXiv licence metadata for this exact version and shown on the abstract page https://arxiv.org/abs/2510.01352v1. A full rights notice is delivered at licenses/arxiv-2510.01352-CC-BY-4.0.txt. Characters: every retained excerpt is pure ASCII, so the delivered engine, XeLaTeX with its default Latin Modern text font, reports no missing character.
\begin{theorem}[Weight formula]\label{thm:weightformula}
  The weight character associated with any descent $(d_1d_2\cdots d_\ell)_n$ is given by,
  \begin{equation*}
    \omega\circ (d_1d_2\cdots d_\ell)_n=(C^n_{d_1}-1)\ast(C^n_{d_2}-1)\ast\cdots\ast(C^n_{d_\ell}-1),
  \end{equation*}
  or equivalently, it is given by,
  \begin{equation*}  
    \omega\circ (d_1d_2\cdots d_\ell)_n=\sum (-1)^{m+\ell}C_{d_{\alpha_1}}^n\ast C_{d_{\alpha_2}}^n\ast\cdots\ast C_{d_{\alpha_m}}^n,
  \end{equation*}
  where the sum is over all ordered subsets $\{\alpha_1,\ldots,\alpha_m\}$ of $\{1,\ldots,\ell\}$ with $m\leqslant\ell$.
\end{theorem}

\begin{lemma}[Sato coefficient endomorphism]\label{lemma:Satocoeffsviadescents}
  The Sato coefficients $b^\ast_{n-k}$ are given by,
  \begin{equation*}
    b^\ast_{n-k}=-\sum_{\|\df\|=k-1} S\circ\df\botimes\df.
  \end{equation*}
\end{lemma}

  \begin{align}
    \sigma_n&\circ S\circ (d_1d_2d_3)_n\nonumber\\
           =&\;-(123-12-13-23+1+2+3-\nu)\nonumber\\
            &\;+(12^=3-12^=-13^=+1)\nonumber\\
            &\;+(123^=-12-23^=+2)\nonumber\\
            &\;+(123-13-23+3)\nonumber\\
            &\;-(12^=3^=-12^=)-(12^=3-13^=)-(123^=-23^=)\nonumber\\
            &\;+12^=3^=.\label{eq:siglength3}
  \end{align}

\begin{remark}\label{rmk:mainthmstrategy}
 We observe that on the right in \eqref{eq:siglength3}, cancellations occur locally between neighbouring terms corresponding to individual terms in $S=-J+J\cendot\id-J\cendot\id^{\cendot2}+J\cendot\id^{\cendot3}$.
 In other words, the terms generated by $-J$, corresponding to the first bracketed term on the right, cancel completely with some terms in $J\cendot\id$, corresponding to the next three bracketed terms on the right.
 And then the remaining terms from $J\cendot\id$ cancel completely with some terms from $-J\cendot\id^{\cendot2}$, corresponding to the next three pairs of bracketed terms on the right.
 Finally, we observe that the  one remaining term in $-J\cendot\id^{\cendot2}$, namely $12^=3^=$, cancels with the final term on the right generated by $J\cendot\id^{\cendot3}$. 
 We exploit this observation in our proof of Theorem~\ref{thm:Satocoeffsdescents}, next.  
\end{remark}

\begin{theorem}[Sato coefficient via descents]\label{thm:Satocoeffsdescents}
  In the case $k=n+1$ in Lemma~\ref{lemma:Satocoeffsviadescents}, the push forward of the Sato coefficient $b^\ast_{-1}$ by $\sigma_n\botimes\id$ is given by,
  \begin{equation}\label{eq:finalSatocoeff}
        b_{-1}=-\sum_{\|\df\|=n}(-1)^{|\df|}\botimes\df,
  \end{equation}
  or equivalently, $\sigma_n\circ S\circ\df=(-1)^{|\df|}$.
\end{theorem}

\begin{proof}
We proceed systematically through the terms generated by each endomorphism in $S$, namely $-J$, $J\cendot\id$, $-J\cendot\id^{\cendot2}$, etc. 
As indicated in Remark~\ref{rmk:mainthmstrategy}, we compare the terms generated by each endomorphism with the next, and keep track of the cancellation of terms.
For convenience here, we use the following notation for ordered subsets.
For a generic ordered subset, say, $\alpha_1\cdots\alpha_m$, we suppose,
\begin{equation*}
\balpha\coloneqq\alpha_1\cdots\alpha_m\qquad\text{and}\qquad\balpha^=\coloneqq\alpha_1^=\alpha_2\cdots\alpha_m.
\end{equation*}
We use $|\balpha|$ to denote the length of the ordered subset, so for example, $|\alpha_1\cdots\alpha_m|=m$.
Recall that $J\circ\emptyset_0=0$ and $\Delta\circ\emptyset_0=0$.
Hereafter, $\df=(d_1d_2\cdots d_\ell)_n$ denotes a generic descent in $\mathbb D$.
To begin, using Theorem~\ref{thm:weightformula}, we know that,
\begin{equation}\label{eq:mainstep1}
  \sigma_n\circ(-J)\circ\df=-\omega\circ\df=-\sum(-1)^{\ell+|\balpha|}\balpha,
\end{equation}
where the sum is over all ordered subsets $\balpha\in[1,\ldots,\ell]$.

Next, we consider the terms generated by $\sigma_n\circ J\cendot\id\circ\df$. We observe that,
\begin{align*}
  \Delta\circ\df=\sum_{k=1}^\ell&\Bigl((d_1\cdots d_{k-1})_{d_k}\\
                               &\!\otimes\!(d_{k+1}-d_k-1)\cdots(d_\ell-d_k-1)_{n-d_k-1}\Bigr),
\end{align*}
where we have ignored the first term starting with a $\emptyset_0$ left factor as this is eliminated in the next step.
Indeed, when computing $\sigma_n\circ J\cendot\id\circ\df$, the first term just mentioned is eliminated, and within the sum above,
we generate the three factors $\hat{\gamma}_n\bigl(d_k(n-d_k-1)\bigr)=C_{d_k}^n$, $\omega\circ(d_1\cdots d_{k-1})_{d_k}$ and $\omega\circ(d_{k+1}-d_k-1)\cdots(d_\ell-d_k-1)_{n-d_k-1}$.
Proceeding as in Step~2 in the proof of Theorem~\ref{thm:weightformula}, we observe that $\sigma_n\circ J\cendot\id\circ\df$ is given by,
\begin{equation}
  \sum_{k=1}^\ell\sum(-1)^{\ell+|\balpha|+|\bbeta|-1}\balpha k\bbeta^=,\label{eq:mainstep2a}
\end{equation}
where the second sum is over all $\balpha\in[1,\ldots,k-1]$ and $\bbeta\in[k+1,\ldots,\ell]$.
Note that the binomial coeeficient corresponding to the first term $\beta_1^=$ in $\bbeta^=$, also has `$d_k$' subtracted from both its arguments.
Consider the case when $\bbeta=\emptyset$, the empty ordered subset. The terms in \eqref{eq:mainstep2a} corresponding to this case are,
\begin{equation}\label{mainstep2b}
  \sum_{k=1}^\ell\sum(-1)^{\ell+|\balpha|-1}\balpha k.
\end{equation}
We note that we have the following natural decomposition,
\begin{equation*}
[1,\ldots,\ell]\backslash\{\emptyset\}=1\cup[1]2\cup[1,2]3\cup\cdots\cup[1,\ldots,\ell-1]\ell.
\end{equation*}
Hence, except for the single term in \eqref{eq:mainstep1} corresponding to $\balpha=\emptyset$, the terms in \eqref{mainstep2b} exactly match those in \eqref{eq:mainstep1}.
With care taken with signs, these terms exactly cancel. We are thus left with the single term from \eqref{eq:mainstep1} mentioned, and the terms from \eqref{eq:mainstep2a}, for which,
\begin{equation}\label{eq:mainconds2}
  \bbeta\in[k+1,\ldots,\ell]\backslash\{\emptyset\}\quad\text{and}\quad k\in1,\ldots,\ell-1,
\end{equation}
where the former necessitates the latter.

We now consider the terms generated by $\sigma_n\circ(-J\cendot\id^{\cendot2})\circ\df$.
To generate $\Delta_2\circ\df$, we simply need to apply $\Delta$ to the left factors in our expression for $\Delta\circ\df$ above.
This eliminates the first term in $\Delta\circ\df$, and means we need to compute $\Delta\circ(d_1\cdots d_{k-1})_{d_k}$ for which can use the formula for $\Delta\circ\df$ itself, replacing $n$ by $d_k$ and $\ell$ by $d_{k-1}$.
Proceeding in this manner, we find that,
\begin{align*}
\Delta_2\circ\df&=\sum_{k_1=1}^\ell\sum_{k_2=1}^{k_1-1}\Bigl((d_1\cdots d_{k_2-1})_{d_{k_2}}\\
&\otimes(d_{k_2+1}-d_{k_2}-1)\cdots(d_{k_1-1}-d_{k_2}-1)_{d_{k_1}-d_{k_2}-1}\\
&\otimes(d_{k_1+1}-d_{k_1}-1)\cdots\cdot(d_\ell-d_{k_1}-1)_{n-d_{k_1}-1}\Bigr).
\end{align*}
where, again, we have ignored all terms starting with a $\emptyset_0$ left factor.
Thus, using $\Delta_2\circ\df$, computing $\sigma_n\circ(-J\cendot\id^{\cendot2})\circ\df$, generates the following four factors within the double sum involving $k_1$ and $k_2$. 
The first factor is,
\begin{equation*}
\hat{\gamma}_n\bigl(d_{k_2}(d_{k_1}-d_{k_2}-1)(n-d_{k_1}-1)\bigr)=C_{d_{k_2}}^nC_{d_{k_1}-d_{k_2}-1}^{n-d_{k_2}-1},
\end{equation*}
while the other three factors are the weights of each of the three tensored terms shown in the double sum for $\Delta_2\circ\df$ shown above.
We distribute the two factors from $\gamma_n$ above across the first two weights.
We observe that, by an exactly analogous calculation to that in Step~2 in the proof of Theorem~\ref{thm:weightformula}, we have,
\begin{equation}\label{eq:mainstep3a}
    C_{d_{k_2}}^n\omega\circ(d_1\cdots d_{k_2-1})_{d_{k_2}}=\sum(-1)^{k_2-1+|\balpha|}\balpha\,k_2.
\end{equation}
The sum is over all ordered subsets $\balpha\in[1,\ldots,k_2-1]$. Next, again by an analogous calculation, we find,
\begin{align}
  C_{d_{k_1}-d_{k_2}-1}^{n-d_{k_2}-1}\omega\!\circ\!(d_{k_2+1}&\!-\!d_{k_2}\!-\!1)\!\cdot\!\cdot\!\cdot\!(d_{k_1-1}\!-\!d_{k_2}\!-\!1)_{d_{k_1}-d_{k_2}-1}\nonumber\\
  =&\sum(-1)^{k_1-k_2-1+|\bbeta|}\bbeta^=\,k_1.\label{eq:mainstep3b}
\end{align}
The sum is over $\bbeta\in[k_2+1,\ldots,k_1-1]$.
Some care is required here, and we remark that the first term $\beta_1^=$ in $\bbeta^=$ is actually, $C_{d_{\beta_1}-d_{k_2}-1}^{n-d_{k_2}-1}$, as we should anticipate.
Lastly, the final weight term is given exactly by the same calculation as that for the corresponding term in Step~2 of the proof of Theorem~\ref{thm:weightformula}, generating,
\begin{equation}\label{eq:mainstep3c}
    \sum(-1)^{\ell-k_1+|\bgamma|}\bgamma^=.
\end{equation}
The sum is over all $\bgamma\in[k_1+1,\ldots,\ell]$, and similar care is required in interpreting the first term in $\bgamma^=$, as we outlined for $\bbeta^=$ just above. 
Putting the factors \eqref{eq:mainstep3a}, \eqref{eq:mainstep3b} and \eqref{eq:mainstep3c} together, we observe that $\sigma_n\circ(-J\cendot\id^{\cendot2})\circ\df$ equals,
\begin{equation}\label{eq:mainstep3}
   -\sum_{k_1=2}^\ell\sum_{k_2=1}^{k_1-1}\sum(-1)^{\ell+|\balpha|+|\bbeta|+|\bgamma|}\balpha\,k_2\bbeta^=\,k_1\bgamma^=.
\end{equation}
The third sum shown is over all $\balpha\in[1,\ldots,k_2-1]$, $\bbeta\in[k_2+1,\ldots,k_1-1]$ and $\bgamma\in[k_1+1,\ldots,\ell]$.
Now consider the case when $\bgamma=\emptyset$, in which case \eqref{eq:mainstep3} collapses to,
\begin{equation*}
   -\sum_{k_1=2}^\ell\sum_{k_2=1}^{k_1-1}\sum(-1)^{\ell+|\balpha|+|\bbeta|}\balpha\,k_2\bbeta^=\,k_1.
\end{equation*}
If we swap the order of the $k_1$ and $k_2$ sums, and then swap the $k_1$ and $k_2$ labels, we obtain,
\begin{equation}\label{eq:mainstep3d}
   -\sum_{k_1=1}^{\ell-1}\sum_{k_2=k_1+1}^{\ell}\sum(-1)^{\ell+|\balpha|+|\bbeta|}\balpha\,k_1\bbeta^=\,k_2,
\end{equation}
where now $\balpha\in[1,\ldots,k_1-1]$ and $\bbeta\in[k_1+1,\ldots,k_2-1]$.
We note that in \eqref{eq:mainstep3d}, the subsum over $k_2$ and $\bbeta$ of $\bbeta^=\,k_2$ equals the sum of $\hat{\bbeta}$
over all $\hat{\bbeta}\in[k_1+1,\ldots,\ell]\backslash\{\emptyset\}$. This is due to the decomposition of $[k+1,\ldots,\ell]\backslash\{\emptyset\}$ into,
\begin{equation}\label{eq:decomp}
(k+1)\cup[k+1](k+2)\cup\cdots\cup[k+1,\ldots,\ell-1]\ell.
\end{equation}
Hence \eqref{eq:mainstep3d} becomes,
\begin{equation}\label{eq:mainstep3e}
   -\sum_{k_1=1}^{\ell-1}\sum(-1)^{\ell+|\balpha|+|\bbeta|-1}\balpha\,k_1\hat{\bbeta}^=,
\end{equation}
where the second sum is over all $\balpha\in[1,\ldots,k_1-1]$ and $\hat{\bbeta}\in[k_1+1,\ldots,\ell]\backslash\{\emptyset\}$.
This exactly matches and cancels with the terms remaining from \eqref{eq:mainstep2a} with the conditions~\eqref{eq:mainconds2}.
We are thus left with the terms in \eqref{eq:mainstep3}, restricted so that,
\begin{equation}\label{eq:mainconds3}
\bgamma\in[k_1+1,\ldots,\ell]\backslash\{\emptyset\}\quad\text{and}\quad k_1\in2,\ldots,\ell-1. 
\end{equation}

As a last case, we consider $\sigma_n\circ(-J\cendot\id^{\cendot3})\circ\df$. The pattern thereafter becomes straightforward. 
To compute $\Delta_3\circ\df$, we must apply $\Delta$ to the left factors in the expression $\Delta_2\circ\df$ above.
Ignoring the term with the left factor $\emptyset_0$, we find,
\begin{align*}
  \Delta_3\circ\df&=\sum_{k_1=1}^\ell\sum_{k_2=1}^{k_1-1}\sum_{k_3=1}^{k_2-1}\Bigl((d_1\cdots d_{k_3-1})_{d_{k_3}}\\
                  &\otimes(d_{k_3+1}-d_{k_3}-1)\cdots(d_{k_2-1}-d_{k_3}-1)_{d_{k_2}-d_{k_3}-1}\\
                  &\otimes(d_{k_2+1}-d_{k_2}-1)\cdots(d_{k_1-1}-d_{k_2}-1)_{d_{k_1}-d_{k_2}-1}\\
                  &\otimes(d_{k_1+1}-d_{k_1}-1)\cdots\cdot(d_\ell-d_{k_1}-1)_{n-d_{k_1}-1}\Bigr).
\end{align*}
Using this form to compute $\sigma_n\circ(-J\cendot\id^{\cendot3})\circ\df$ generates five factors within the triple sum,
the first of which is the $\hat{\gamma}_n$ factor, which in this case is,
\begin{equation}
  C_{d_{k_3}}^nC_{d_{k_2}-d_{k_3}-1}^{n-d_{k_3}-1}C_{d_{k_1}-d_{k_2}-1}^{n-d_{k_2}-1}. \label{eq:threefactors}
\end{equation}
The other four factors are the weights of the four tensored terms shown in the triple sum.
We distribute the three binomial coefficients respectively across the first three weights.
Again, exactly analogous to the calculation in Step~2 in the proof of Theorem~\ref{thm:weightformula}, we have,
\begin{equation}\label{eq:mainstep4a}
    C_{d_{k_3}}^n\omega\circ(d_1\cdots d_{k_3-1})_{d_{k_3}}=\sum(-1)^{k_3-1+|\balpha|}\balpha\,k_3.
\end{equation}
The sum is over all $\balpha\in[1,\ldots,k_3-1]$.
The next two factors are computed in the same manner as that above. Indeed, respectively,
the second/third factor in \eqref{eq:threefactors} times the weight of the second/third factor in $\Delta_3\circ\df$ give,
\begin{align}
  &\sum(-1)^{k_2-k_3-1+|\bbeta|}\bbeta^=\,k_2, \label{eq:mainstep4b}
  \intertext{and}
  &\sum(-1)^{k_1-k_2-1+|\bgamma|}\bgamma^=k_1. \label{eq:mainstep4c}
\end{align}
The sums are respectively over $\bbeta\in[k_3+1,\ldots,k_2-1]$ and $\bgamma\in[k_2+1,\ldots,k_1-1]$.
The final weight term generates,
\begin{equation}\label{eq:mainstep4d}
    \sum(-1)^{\ell-k_1+|\bdelta|}\bdelta^=.
\end{equation}
The sum is over $\bdelta\in[k_1+1,\ldots,\ell]$.
Putting the factors \eqref{eq:mainstep4a}--\eqref{eq:mainstep4d} together, we observe, $\sigma_n\circ(-J\cendot\id^{\cendot3})\circ\df$ equals,
\begin{equation}\label{eq:mainstep4}
   \sum_{k_1=3}^\ell\sum_{k_2=2}^{k_1-1}\sum_{k_3=1}^{k_2-1} \sum\rho\,\balpha\,k_3\bbeta^=\,k_2\bgamma^=k_1\bdelta^=,
\end{equation}
where $\rho=(-1)^{\ell+|\balpha|+|\bbeta|+|\bgamma|+|\bdelta|-1}$.
The fourth sum is over all ordered subsets associated with the factors \eqref{eq:mainstep4a}--\eqref{eq:mainstep4d}.
Now consider the case when $\bdelta=\emptyset$, in which case \eqref{eq:mainstep4} collapses to, 
\begin{equation}\label{eq:mainstep4e}
   \sum_{k_1=3}^\ell\sum_{k_2=2}^{k_1-1}\sum_{k_3=1}^{k_2-1} \sum(-1)^{\ell+|\balpha|+|\bbeta|+|\bgamma|-1}\balpha\,k_3\bbeta^=\,k_2\bgamma^=k_1,
\end{equation}
where the fourth sum is now only over the ordered subsets for $\balpha$, $\bbeta$ and $\bdelta$.
If we reverse the order of the $k_1$, $k_2$ and $k_3$, which involves three swaps, and then swap the $k_1$ and $k_3$ labels, we obtain,
\begin{equation}\label{eq:mainstep4f}
   \sum_{k_1=1}^{\ell-2}\sum_{k_2=k_1+1}^{\ell-1}\sum_{k_3=k_2+1}^{\ell} \sum\rho^\prime\balpha\,k_1\bbeta^=\,k_2\bgamma^=k_3,
\end{equation}
where $\rho^\prime=(-1)^{\ell+|\balpha|+|\bbeta|+|\bgamma|-1}$.
The fourth sum is over $\balpha\in[1,\ldots,k_1-1]$, $\bbeta\in[k_1+1,\ldots,k_2-1]$ and $\bgamma\in[k_2+1,\ldots,k_3-1]$.
Again using the decomposition \eqref{eq:decomp}, the form \eqref{eq:mainstep4f} collapses to,
\begin{equation}\label{eq:mainstep4g}
   \sum_{k_1=1}^{\ell-2}\sum_{k_2=k_1+1}^{\ell-1}\sum(-1)^{\ell+|\balpha|+|\bbeta|+|\bgamma|}\balpha\,k_1\bbeta^=\,k_2\hat{\bgamma}^=,
\end{equation}
where the third sum is over the same ranges for $\balpha$ and $\bbeta$, but now $\hat{\bgamma}\in[k_2+1,\ldots,\ell]\backslash\{\emptyset\}$.
If we swap the order of the $k_1$ and $k_2$ sums, and then swap the $k_1$ and $k_2$ labels, in \eqref{eq:mainstep4g},
then it precisely matches \eqref{eq:mainstep3} subject to the conditions \eqref{eq:mainconds3}, with opposite sign.
These two terms thus cancel.
The procedure is now straightforward. We continue in this manner until we reach the final single term $12^=3^=\cdots\ell^=$, which cancels with the corresponding term at the previous level.
There is thus only a single term remaining, namely the `$-(-1)^{\ell}$' term from \eqref{eq:mainstep1} corresponding to the case $\balpha=\emptyset$.
This completes the proof.
\qed
\end{proof}

\end{document}
